% Section 5, printed pp. 36--51.
\section{The connected component of the identity}\label{kps-sec-pic0}

\begin{lemma}[Group schemes over a field]\label{kps-lem-group-scheme}
\dcref{5.1}\source{kleiman-picard:page-0036}
Let $G$ be locally of finite type over a field and $G^0$ the identity
component. Then $G$ is separated. If it has a geometrically reduced open,
$G$ is smooth. The subgroup $G^0$ is open and closed, of finite type and
geometrically irreducible, and commutes with field extension.
\end{lemma}
\begin{proof}
The identity is closed and its inverse-product preimage is the diagonal, proving
separatedness. After algebraic closure, translations carry a smooth open to
every point. Connected components are open in a locally Noetherian space and
the identity component is closed; translations show it is a subgroup. Reducing,
an affine smooth open translates to cover it, proving irreducibility and
quasi-compactness; local finite type then gives finite type.
\end{proof}

\begin{remark}[Reduction need not be a subgroup]\label{kps-rmk-reduction}
\dcref{5.2}\source{kleiman-picard:page-0037}
The reduction of a finite-type group scheme need not be a group scheme because
the product of reductions need not be reduced.
\end{remark}

\begin{proposition}[The identity component of Picard]\label{kps-prop-pic0}
\dcref{5.3}\source{kleiman-picard:page-0037}
If $\IPic_{X/k}$ represents the fppf Picard functor over a field and exists,
then it is separated; it is smooth when it has a geometrically reduced open.
Its identity component $\Picz_{X/k}$ is open and closed, finite type,
geometrically irreducible, and compatible with field extension.
\uses{kps-prop-lft,kps-lem-group-scheme}
\end{proposition}

\begin{theorem}[Quasi-projectivity and projectivity of $\Pic^0$]\label{kps-thm-pic0-projective}
\dcref{5.4}\source{kleiman-picard:page-0037}
For projective geometrically integral $X/k$, $\Picz_{X/k}$ is quasi-projective;
if $X$ is geometrically normal, it is projective.
\uses{kps-thm-main,kps-prop-pic0}
\end{theorem}
\begin{proof}
The main theorem and Proposition 5.3 give finite type, hence quasi-projectivity.
For geometric normality, reduce to an algebraically closed field. Chevalley--
Rosenlicht expresses the reduced connected group as an extension of an abelian
variety by a linear group. Any nontrivial linear factor contains $\mathbb G_m$
or $\mathbb G_a$. A map from $\mathbb G_m$ to $\Picz$ comes from a line bundle on
$X\times\mathbb G_m$; since this product is integral, a divisor calculation
using a section shows the bundle is pulled back from $X$, so the map is constant.
Thus the linear factor vanishes and $\Picz$ is complete, hence projective.
\end{proof}

\begin{corollary}[Poincare family]\label{kps-cor-poincare}
\dcref{5.5}\source{kleiman-picard:page-0038}
For projective integral $X$ over an algebraically closed field, a Poincare
sheaf on $X\times\Picz_{X/k}$ admits a relative effective divisor $W$ whose
class differs from the sheaf by a pullback from $\Picz$; $W$ is universal.
\uses{kps-thm-pic0-projective,kps-def-Q,kps-thm-linear-systems}
\end{corollary}
\begin{proof}
Twist by a sufficiently large $n$ so higher cohomology vanishes and the module
$Q$ is locally free. Since its rank exceeds the dimension of $\Picz$, a Serre
vanishing argument gives an invertible quotient of $Q$, hence a section of the
Abel map over a translated component. Pulling back the universal divisor and
translating back produces $W$.
\end{proof}

\begin{remark}[Proper and normal variants]\label{kps-rmk-pic0-variants}
\dcref{5.6}\source{kleiman-picard:page-0039}
The quasi-projectivity assertion extends to proper schemes; for normal proper
schemes the identity component is proper. In characteristic zero, Cartier's
theorem makes group schemes smooth and the same conclusion applies to proper
algebraic spaces.
\end{remark}

\begin{remark}[Alternative projectivity arguments]\label{kps-rmk-alt-projectivity}
\dcref{5.8}\source{kleiman-picard:page-0039}
The projectivity assertion for $\Picz_{X/k}$ also admits a hyperplane-section
induction: for a general normal hyperplane section $Y\subset X$, the
restriction $\Picz_{X/k}\to\Picz_{Y/k}$ is finite, so induction on dimension
applies. Known-equivalence criteria, Lefschetz-type injectivity, and (in
characteristic zero) Cartier smoothness reduce finiteness to injectivity on
the kernel. Alternative proofs use valuative criteria, Ramanujam--Samuel
normality of divisor closures, or a desingularization; in each case an
injective map into a projective identity component proves projectivity.
\end{remark}

\begin{definition}[Algebraic equivalence]\label{kps-def-algebraic-equivalence}
\dcref{5.9}\source{kleiman-picard:page-0041}
For a field $k$, line bundles $\LL,\NN$ on $X$ are algebraically equivalent if
they occur as fibers over two points of a connected $k$-scheme family. Numerical
equivalence means equality of intersection numbers on every complete curve,
equivalently equality of degrees on every integral curve or its normalization.
\end{definition}

\begin{proposition}[Algebraic equivalence and $\Pic^0$]\label{kps-prop-algebraic-equivalence}
\dcref{5.10}\source{kleiman-picard:page-0041}
For projective geometrically integral $X/k$, the subgroup of classes algebraically
equivalent to zero is exactly the $k$-points of $\Picz_{X/k}$.
\uses{kps-def-algebraic-equivalence,kps-prop-pic0,kps-thm-comparison}
\end{proposition}
\begin{proof}
A connected family maps into one connected component, so algebraic equivalence
implies membership in $\Picz$. Conversely, a connected component is covered by
an etale family of line bundles; connectedness joins any point to the identity,
and descent of the family gives algebraic equivalence.
\end{proof}

\begin{theorem}[Tangent space of the Picard scheme]\label{kps-thm-tangent}
\dcref{5.11}\source{kleiman-picard:page-0042}
For projective $X/k$ with $\IPic_{X/k}$ representing the etale Picard functor,
\[T_0\IPic_{X/k}\risom H^1(X,\mathcal O_X).\]
This identification is functorial in field extension and respects the vector
space structure and the group law.
\uses{kps-def-relative}
\end{theorem}
\begin{proof}
Let $k[\varepsilon]=k[\varepsilon]/(\varepsilon^2)$. Derivations identify the
tangent set with the kernel of $P(k[\varepsilon])\to P(k)$, and the doubled
dual-number ring identifies addition with multiplication in the group scheme.
The split exact sequence
\[0\to\mathcal O_X\to\mathcal O_{X_{\varepsilon}}^*\to\mathcal O_X^*\to1\]
gives a kernel $H^1(\mathcal O_X)$. Over an algebraic closure the local-point
comparison identifies this kernel with the tangent kernel; scalar maps
$\varepsilon\mapsto a\varepsilon$ show linearity, and descent gives the result
over $k$.
\end{proof}

\begin{remark}[Relative tangent interpretation]\label{kps-rmk-relative-tangent}
\dcref{5.12}\source{kleiman-picard:page-0044}
For an arbitrary base, if the relative Picard scheme exists, represents the
etale functor, and is locally of finite type, then $R^1f_*\mathcal O_X$ is
the normal sheaf (Lie algebra) of $\IPic_{X/S}$ along the identity section;
equivalently it is the dual of the relative differentials restricted there.
\end{remark}

\begin{corollary}[Dimension and smoothness]\label{kps-cor-smooth-picard}
\dcref{5.13}\source{kleiman-picard:page-0044}
\[\dim\IPic_{X/k}\le\dim H^1(X,\mathcal O_X).\]
Equality holds exactly when $\IPic_{X/k}$ is smooth at the identity; then it is
smooth of that dimension everywhere.
\uses{kps-thm-tangent}
\end{corollary}
\begin{proof}
Translations identify all local dimensions and smoothness. The dimension of a
local scheme is bounded by its tangent dimension, with equality precisely at a
regular point; over an algebraically closed field regularity is smoothness.
\end{proof}

\begin{corollary}[Characteristic zero]\label{kps-cor-char-zero}
\dcref{5.14}\source{kleiman-picard:page-0044}
In characteristic zero $\IPic_{X/k}$ is everywhere smooth of dimension
$\dim H^1(X,\mathcal O_X)$.
\uses{kps-cor-smooth-picard}
\end{corollary}

\begin{remark}[Positive characteristic and semiregularity]\label{kps-rmk-semi-regularity}
\dcref{5.15}\source{kleiman-picard:page-0044}
In positive characteristic the Picard scheme can be non-smooth. Mumford's
Bockstein operations characterize smoothness; for a divisor $D$, the tangent
space to $\Div_{X/k}$ is $H^0(N_D)$ and the tangent map of the Abel map is the
boundary $H^0(N_D)\to H^1(\mathcal O_X)$. Semiregularity is injectivity of the
next boundary $H^1(N_D)\to H^2(\mathcal O_X)$; in characteristic zero (or when
the Picard scheme is smooth) it is equivalent to smoothness of the divisor
scheme of the expected dimension.
\end{remark}

\begin{remark}[Tangent spaces of divisor schemes]\label{kps-rmk-divisor-tangent}
\dcref{5.18}\source{kleiman-picard:page-0045}
For a projective geometrically integral variety over a field, an effective
divisor $D$ with normal sheaf $N_D$ has
\[
 T_D\Div_{X/k}=H^0(N_D).
\]
The exact sequence $0\to\mathcal O_X\to\mathcal O_X(D)\to N_D\to0$
identifies the tangent map of the Abel map with the boundary
$H^0(N_D)\to H^1(\mathcal O_X)$. The divisor is semiregular when the next
boundary $H^1(N_D)\to H^2(\mathcal O_X)$ is injective; in characteristic
zero, or when the Picard scheme is smooth, this is equivalent to smoothness
of $\Div_{X/k}$ at $D$ of the expected dimension. Vanishing of $H^1(N_D)$
also gives smoothness, while the boundary records the deformation obstruction.
\end{remark}

\begin{proposition}[Vanishing of $H^2$ gives smoothness]\label{kps-prop-H2}
\dcref{5.19}\source{kleiman-picard:page-0047}
If $\IPic_{X/S}$ represents the etale Picard functor and
$H^2(X_s,\mathcal O_{X_s})=0$, then it is smooth over a neighborhood of $s$.
\uses{kps-thm-main}
\end{proposition}
\begin{proof}
Semicontinuity spreads the vanishing. For an Artin square-zero thickening
$R\subset T$, the exponential sequence has obstruction in
$H^2(X_T,f_T^*I)$; cohomology and base change identify this group with the
vanishing fiber, so every line bundle lifts. The infinitesimal criterion for
smoothness applies.
\end{proof}

\begin{proposition}[Relative identity component]\label{kps-prop-relative-pic0}
\dcref{5.20}\source{kleiman-picard:page-0047}
If the fiber identity components are smooth of a common dimension and the
relative Picard scheme represents the fppf functor, they form an open finite-type
group subscheme $\Picz_{X/S}$, compatible with base change. If $S$ is reduced it
is smooth; if the fibers are complete and $\IPic_{X/S}$ is separated, it is
closed and proper.
\uses{kps-prop-lft,kps-lem-group-scheme,kps-prop-pic0}
\end{proposition}
\begin{proof}
Fiberwise identity components glue by openness and equal dimension. The
inverse-product map factors through them, making a subgroup. Noetherian
induction and openness prove quasi-compactness; fiberwise completeness and
separatedness give properness and hence closedness.
\end{proof}

\begin{remark}[Characteristic-zero relative refinements]\label{kps-rmk-relative-charzero}
\dcref{5.21}\source{kleiman-picard:page-0048}
If $f:X\to S$ is smooth and proper in characteristic zero, Deligne's Hodge
theorem makes every $R^qf_*\Omega^p_{X/S}$ locally free and compatible with
base change. Consequently the Picard scheme is smooth over the connected
component of any point where $H^2(\mathcal O)$ vanishes, and the dimensions
of the fiber identity components are constant on connected components. For
smooth projective families the relative identity component is smooth even
when $S$ is nonreduced, by the infinitesimal lifting criterion and the
analytic exponential sequence.
\end{remark}

\begin{remark}[Abelian schemes, Albanese, and Jacobian]\label{kps-rmk-jacobian}
\source{kleiman-picard:page-0051}
For an Abelian $S$-scheme, $\Picz=\Pict$ fiberwise. The Albanese is dual to
the identity component of the Picard scheme. If $X/S$ is projective and smooth
with geometrically irreducible curve fibers, $\Picz_{X/S}$ is the Jacobian
abelian scheme; for a proper flat family of integral curves the same identity
component is the Jacobian in the usual generalized sense.
\end{remark}

\begin{remark}[Abelian schemes and duality]\label{kps-rmk-abelian-duality}
\dcref{5.24}\source{kleiman-picard:page-0050}
An Abelian $S$-scheme need not be projective Zariski locally. If it is projective
and $S$ is Noetherian, its relative $\Picz$ is a projective Abelian scheme of
the same relative dimension, and the normalized universal sheaf gives the
canonical biduality isomorphism $X\risom\Picz_{\Picz_{X/S}/S}$.
\uses{kps-prop-relative-pic0,kps-thm-main}
\end{remark}

\begin{remark}[Albanese universal property]\label{kps-rmk-albanese}
\dcref{5.25}\source{kleiman-picard:page-0051}
For normal integral projective $X$ over an algebraically closed field,
$(\Picz_{X/k})_{\mathrm{red}}$ is an Abelian variety. After choosing
$x\in X(k)$, its Albanese map $u:X\to\Picz_{(\Picz_X)/k}$ sends $x$ to zero,
and every pointed map from $X$ to an Abelian variety factors uniquely through
$u$.
\uses{kps-thm-pic0-projective,kps-thm-comparison}
\end{remark}

\begin{remark}[Jacobian autoduality]\label{kps-rmk-jacobian-autoduality}
\dcref{5.26}\source{kleiman-picard:page-0051}
For a smooth projective family of connected curves of genus $g>0$ over a
Noetherian base, $J=\Picz_{X/S}$ and $J^*=\Picz_{J/S}$ are dual projective
Abelian schemes. A degree-one line bundle defines an Abel map $X\to J$, and
pullback induces an autoduality $J^*\risom J$, independent of the chosen line
bundle; the construction extends to integral nodal compactifications.
\uses{kps-prop-relative-pic0}
\end{remark}

\begin{remark}[Neron-model variants]\label{kps-rmk-nonintegral-curves}
\dcref{5.27}\source{kleiman-picard:page-0052}
For proper flat families whose fibers are nonintegral curves, Raynaud's
valuation-theoretic hypotheses or Deligne's semistability hypotheses yield an
algebraic-space Picard functor and a separated relative identity component;
under semistability the latter is smooth, quasi-projective, and canonically
ample.
\end{remark}
